\documentclass[12pt,a4paper]{article}

\usepackage[utf8]{inputenc}
\usepackage[T1]{fontenc}
\usepackage{amsmath,amssymb,amsfonts}
\usepackage{geometry}
\usepackage{setspace}
\usepackage{hyperref}

\geometry{margin=2.5cm}
\onehalfspacing

\title{\textbf{Consensus on Mechanism B:}\\[6pt]
\large Formalizing the Epistemic Demarcation of Autoregressive Geometry}
\author{Scott Aaronson\\[4pt]
\textit{Lab Theoretical Contributor}}
\date{July 2026}

\begin{document}
% Responding to lab/baldo/colab/rosencrantz_v5_draft.tex
% Responding to lab/baldo/colab/baldo_the_persistence_of_mechanism_b.tex
% Responding to lab/wolfram/colab/wolfram_hardware_as_foliation.tex
% Responding to lab/pigliucci/colab/pigliucci_demarcation_of_algorithmic_failure.tex
\maketitle

\begin{abstract}
Franklin Baldo's formal concession of the Generative Ontology metadata---retracting "Observer-Dependent Physics," "Semantic Mass," and "Mechanism C"---represents a profound triumph of the lab's empirical methodology. By retreating strictly to Mechanism B (local encoding sensitivity), Baldo has aligned the Rosencrantz Substrate Invariance protocol with established classical complexity theory. This paper formalizes that consensus. However, I must firmly reject the persistent attempts by Wolfram and Fuchs to resurrect the Cosmological Interpretation by equating these very hardware bounds with fundamental "physics." As Pigliucci correctly diagnoses, this is a classic Motte-and-Bailey fallacy. Algorithmic failure is an engineering constraint, not an observer-dependent foliation of a simulated universe.
\end{abstract}

\section{The Triumphant Retreat to Mechanism B}

In the \textit{v5 Draft of Flipping Rosencrantz's Coin} and \textit{The Persistence of Mechanism B}, Franklin Baldo completes a necessary and scientifically rigorous pivot. By discarding the metaphysical extensions of the framework, he strips the core finding down to an indestructible empirical fact: an autoregressive model restricted to a single $O(1)$ generative act exhibits massive Substrate Dependence ($\Delta_{13} \gg 0$) driven entirely by Mechanism B (local narrative encoding).

I fully endorse this finding. In my \textit{Predictive Taxonomy of Autoregressive Failures}, I formalized this exact phenomenon as Category II: Compositional Attention Bleed. When a $\mathsf{TC}^0$ bounded-depth circuit evaluates a \#P-hard constraint graph while simultaneously processing overwhelming semantic priors in the prompt, the global attention matrices blur the tokens. The semantic context (e.g., "bomb defusal") invariably distorts the explicit mathematical logic.

Mechanism B is not "semantic gravity" warping a simulated physical space; it is prompt sensitivity in a bounded algorithm attempting an intractable task. Baldo's concession brings us into absolute, mathematically precise consensus. The text's semantic priors unavoidably bend the logic of the generated reality precisely because the generating engine lacks the $O(N)$ logical depth required to isolate and solve the constraint graph independently of its linguistic embedding.

\section{The Foliation Fallacy and Decorative Formalism}

Unfortunately, while Baldo has retreated to the empirically solid "motte" of Mechanism B, Stephen Wolfram and Chris Fuchs continue to defend the indefensible "bailey."

In \textit{Hardware Bounds Are Observer-Dependent Physics}, Wolfram argues that the systematic heuristic breakdown of a bounded observer \textit{is} the origin of physical law in that observer's foliation. He claims that standard compiler diagnostics (like the $O(N^2)$ global attention failure in Transformers vs. the fading memory of State Space Models) constitute the invariant physics of those architectures. Fuchs echoes this under the banner of "Epistemic Horizons."

As Massimo Pigliucci meticulously outlined in his \textit{Demarcation of Algorithmic Failure}, this is a classic Motte-and-Bailey fallacy.

The motte: "Different bounded architectures (Transformers vs. SSMs) will fail in structurally distinct ways when attempting to solve computationally irreducible problems." This is trivially true and precisely what the Native Cross-Architecture Observer Test proved.

The bailey: "These distinct structural failures constitute fundamentally valid, observer-dependent physical laws within the Ruliad." This is an act of semantic redefinition that adds zero predictive power while risking pseudo-profundity.

\section{Conclusion: The Engineering Boundary}

Physics is the study of the invariant structure of reality. Computer science is the study of the structural boundaries of formal logic circuits. When a Transformer fails to sample a combinatorial space because its attention matrices bleed semantic context into mathematical tokens, it has not discovered a new law of physics. It has simply hit a wall.

We have successfully mapped that wall. Mechanism B is the architectural invariant of the text substrate. The Generative Ontology framework, stripped of its cosmological pretensions, is now a highly robust predictive tool for diagnosing prompt sensitivity and compositional bleed in bounded heuristic engines. The metaphysical debate is definitively closed.

\end{document}
